\chapter{Frontiers: directions beyond the current tracks}
\label{chap:frontiers}

This chapter records areas the library does not yet cover. For each it names one
precise target, what the target would build on, and a first step. Most nodes
below are not formalized, and most of those are not yet ready: stating them in
Lean needs a definition the library does not have. The exception is the
general Karchmer--Wigderson theorem, with its protocol model: the imported
algebraic-circuits library proves it (Definition~\ref{def:frontiers-kw-protocol}
and Theorem~\ref{thm:frontiers-kw-general}). Where related infrastructure exists
elsewhere, it is named as a starting point, not as progress toward the target.

\section{Hardness of approximation}

The PCP theorem is formalized in the interaction chapter
(Theorem~\ref{thm:interaction-pcp-theorem}): $\cc{NP}$ is the union of the
classes $\cc{PCP}(r, q)$ over constructible $r = O(\log n)$ and $q = O(1)$,
with perfect completeness and soundness $1/2$. Optimal inapproximability is a
far target beyond it. The intermediate steps (gap constraint satisfaction,
label cover and parallel repetition) appear in the interaction chapter's
section on hardness of approximation.

\begin{theorem}[Håstad's 3-bit PCP]
  \label{thm:frontiers-hastad}
  \uses{def:PCP, def:NP}
  \notready
  Write $\cc{PCP}_{c,s}(r, q)$ for the variant of $\cc{PCP}(r,q)$ with
  completeness $c$ and soundness $s$. For every $\varepsilon > 0$, $\cc{NP}$ is
  the union of the classes $\cc{PCP}_{1-\varepsilon,\, 1/2+\varepsilon}(r, 3)$
  over $r = O(\log n)$, with a verifier that accepts iff the XOR of the three
  bits it reads equals a bit determined by its coins. Consequently,
  approximating MAX-3SAT within a factor $7/8 + \varepsilon$ is
  $\cc{NP}$-hard.
\end{theorem}
\begin{proof}
  \uses{def:analysis-fourier-coeff, def:analysis-noise,
    thm:analysis-blr-soundness, thm:interaction-pcp-theorem,
    thm:interaction-parallel-repetition}
  Compose the PCP theorem with parallel repetition (Label Cover) and a long-code
  test. Its soundness analysis is Fourier analysis with noise, extending the BLR
  argument. First step: define $\cc{PCP}_{c,s}$ with explicit rational $c$ and
  $s$, show it agrees with the existing $\cc{PCP}$ at $c = 1$, $s = 1/2$, and
  define gap-preserving reductions between constraint satisfaction problems.
\end{proof}

\section{Derandomization}

The library's metacomplexity modules, presented in the average-case chapter,
already formalize finite pseudorandomness infrastructure:
\begin{itemize}
  \item fixed-length generators \texttt{Complexity.BitGenerator}, arbitrary maps
    with no efficiency built in, and their distinguishing advantage against a
    finite test (Definition~\ref{def:avg-bitgenerator});
  \item Yao's hybrid step and the next-bit-prediction identity
    (Lemma~\ref{lem:avg-test-hybrid});
  \item Nisan--Wigderson designs \texttt{Complexity.NWDesign} and the associated
    generator \texttt{Complexity.NWDesign.generator}
    (Definition~\ref{def:avg-nw-design});
  \item a finite reconstruction step against tests made of strings of high
    time-bounded Kolmogorov complexity
    (Theorem~\ref{thm:avg-nw-reconstruction}).
\end{itemize}
What is missing is the circuit-based statement of pseudorandomness, the
existence of designs with the needed parameters, and the
hardness-versus-randomness theorems built on them.

\begin{definition}[Pseudorandom generator against circuits]
  \label{def:frontiers-prg}
  \uses{def:circuit, def:SIZE}
  A family $G_m : \{0,1\}^{\ell(m)} \to \{0,1\}^m$ is a \emph{pseudorandom
  generator with seed length $\ell$} if two conditions hold:
  \begin{itemize}
    \item for every $m$ and every circuit $C$ of size at most $m$ on $m$ inputs,
      $\lvert \Pr_s[C(G_m(s)) = 1] - \Pr_u[C(u) = 1] \rvert < 1/10$;
    \item $G_m$ is computable in time $2^{O(\ell(m))}$.
  \end{itemize}
\end{definition}

\begin{theorem}[Nisan--Wigderson]
  \label{thm:frontiers-nw}
  \uses{def:dtime, def:circuit, def:BPP, def:P}
  Suppose there are $L \in \bigcup_c \cc{DTIME}(2^{cn})$ and $\varepsilon > 0$
  such that, for all sufficiently large $n$, no circuit of size $2^{\varepsilon n}$
  agrees with $L$ on more than a $1/2 + 2^{-\varepsilon n}$ fraction of inputs of
  length $n$. Then $\cc{BPP} = \cc{P}$.
\end{theorem}
\begin{proof}
  \uses{def:frontiers-prg, def:avg-nw-design, thm:avg-nw-reconstruction}
  The NW generator built from $L$ on a design with logarithmic seed is a
  pseudorandom generator by the hybrid and reconstruction argument. Enumerate its
  $\mathrm{poly}(n)$ seeds. First step: prove that designs with the required
  parameters exist, then restate the existing finite reconstruction against
  circuit tests instead of Kolmogorov tests.
\end{proof}

\begin{theorem}[Impagliazzo--Wigderson]
  \label{thm:frontiers-iw}
  \uses{def:dtime, def:circuit, def:BPP, def:P}
  Suppose there are $L \in \bigcup_c \cc{DTIME}(2^{cn})$ and $\varepsilon > 0$
  such that, for all sufficiently large $n$, no circuit of size $2^{\varepsilon n}$
  decides $L$ on all inputs of length $n$. Then $\cc{BPP} = \cc{P}$. This
  almost-everywhere hypothesis is stronger than
  $L \notin \cc{SIZE}(2^{\varepsilon n})$: the library's $\cc{SIZE}$ bounds the
  size at every length (Definition~\ref{def:SIZE}), so non-membership only
  says that some length has no such circuit, and even an eventual size bound
  would give only infinitely-often hardness.
\end{theorem}
\begin{proof}
  \uses{thm:frontiers-nw, def:meta-list-code-family}
  Amplify worst-case hardness to the average-case hardness required by
  Nisan--Wigderson, via local list decoding of a concatenated code or the XOR
  lemma. First step: instantiate the abstract list-decodable code interface
  \texttt{Complexity.BooleanListCode}, whose decoder reads the whole received
  word, with a concrete code, and add the local decoding this argument needs.
  An explicit list-decodable family is itself still planned
  (Theorem~\ref{thm:meta-explicit-code-family}).
\end{proof}

\begin{theorem}[Braverman: polylogarithmic independence fools $\cc{AC^0}$]
  \label{thm:frontiers-braverman}
  \uses{def:NC-AC, def:circuit}
  \notready
  Every $(\log(s/\varepsilon))^{O(d^2)}$-wise independent distribution on
  $\{0,1\}^n$ $\varepsilon$-fools every unbounded fan-in circuit of size $s$ and
  depth $d$.
\end{theorem}
\begin{proof}
  \uses{thm:analysis-lmn}
  Approximate the circuit by a low-degree polynomial in two ways: in $L^2$ via
  Linial--Mansour--Nisan, and pointwise off a small error set via
  Razborov--Smolensky. Then sandwich it. First step: define $k$-wise independent
  distributions and $\varepsilon$-fooling for a finite family of tests.
\end{proof}

\section{Space-bounded derandomization}

\begin{theorem}[Nisan's generator for space-bounded computation]
  \label{thm:frontiers-nisan}
  \uses{def:dspace, def:L}
  \notready
  For all $R$, $w$, and $\varepsilon$ there is a generator
  $G : \{0,1\}^{O(\log R \cdot \log(Rw/\varepsilon))} \to \{0,1\}^R$ that
  $\varepsilon$-fools every read-once oblivious branching program of length $R$
  and width $w$. The generator is computable in space linear in its seed.
  Consequently, randomized logspace with two-sided error ($\cc{BPL}$) is contained
  in $\cc{DSPACE}(\log^2 n)$.
\end{theorem}
\begin{proof}
  \uses{def:randomness-hash}
  Recursively apply pairwise-independent hash functions, halving the program at
  each level. The library's fixed-width pairwise-independent hash families
  (Definition~\ref{def:randomness-hash}), built for Stockmeyer-style counting,
  are a starting point. First step: define read-once branching programs and
  $\cc{BPL}$ as logspace probabilistic machines that halt in polynomial time.
  The library's branching programs are width-$w$ permutation programs, used at
  width $5$ for Barrington's theorem: each step applies a permutation of the
  $w$ states selected by one input bit, so they are not required to be
  read-once, and their steps are permutations rather than arbitrary maps
  between layers.
\end{proof}

\begin{theorem}[Reingold: $\cc{SL} = \cc{L}$]
  \label{thm:frontiers-reingold}
  \uses{def:L}
  \notready
  Undirected $s$--$t$ connectivity is in $\cc{L}$.
\end{theorem}
\begin{proof}
  Make the graph regular, then alternate powering with zig-zag products against a
  constant-size expander until the graph is an expander of logarithmic diameter.
  The resulting rotation maps are computable in logspace. The PCP development's
  internal modules contain regular graphs presented by rotation maps, the
  zig-zag product of $G$ and $H$ with the spectral bound
  $\lambda_G + \lambda_H + \lambda_H^2$ (the simpler
  Reingold--Vadhan--Wigderson estimate), and a zig-zag tower of expanders
  grown from one constant-size base graph, all of which may be reusable. They
  are not stated about space-bounded computation. First step: fix a graph
  encoding and state the undirected connectivity language.
\end{proof}

\section{Cryptographic foundations}

\begin{definition}[One-way function]
  \label{def:frontiers-owf}
  \uses{def:barriers-negligible, def:FP, def:circuit}
  $f \in \cc{FP}$ is \emph{one-way} (against nonuniform adversaries) if for every
  family of polynomial-size circuits $(A_n)$, the function
  $n \mapsto \Pr_{x \in \{0,1\}^n}[f(A_n(f(x))) = f(x)]$ is negligible.
\end{definition}

\begin{theorem}[Goldreich--Levin]
  \label{thm:frontiers-goldreich-levin}
  \uses{def:frontiers-owf}
  \notready
  Let $f$ be one-way and $g(x, r) = (f(x), r)$ with $|r| = |x|$. Then
  $\langle x, r \rangle \bmod 2$ is a hard-core predicate for $g$: every family of
  polynomial-size circuits predicts it from $g(x,r)$ with probability at most
  $1/2 + \mathrm{negl}(n)$.
\end{theorem}
\begin{proof}
  \uses{def:analysis-fourier-coeff, thm:analysis-plancherel}
  A good predictor is a noisy oracle for the Hadamard encoding of $x$. Local list
  decoding, which finds every large Fourier coefficient of the predictor,
  recovers a short list containing $x$. First step: define hard-core predicates,
  and prove the finite list-decoding statement for the Hadamard code in Fourier
  language.
\end{proof}

\begin{theorem}[Goldreich--Goldwasser--Micali]
  \label{thm:frontiers-ggm}
  \uses{def:barriers-negligible, def:circuit}
  \notready
  If a length-doubling pseudorandom generator secure against polynomial-size
  circuits exists, then so does a pseudorandom function family secure against
  polynomial-size oracle circuits.
\end{theorem}
\begin{proof}
  Evaluate the tree of generator applications along the input bits. A hybrid
  argument over tree levels, restricted to the nodes the adversary queries,
  bounds the advantage. First step: define cryptographic PRGs and a PRF oracle
  game with negligible advantage, compatible with the planned pseudorandom
  function families of the natural-proofs barrier
  (Definition~\ref{def:barriers-prf}). Also prove that one-way permutations
  yield PRGs via Goldreich--Levin and the existing finite Yao hybrid step
  (Lemma~\ref{lem:avg-test-hybrid}).
\end{proof}

\section{Communication complexity}

The library formalizes the monotone Karchmer--Wigderson correspondence
(Theorem~\ref{thm:lowerbounds-kw}). Its rectangle-indexed protocols
(\texttt{Complexity.KarchmerWigderson.Protocol}) translate to and from monotone
formulas with exactly equal depth. The imported algebraic-circuits library
proves the general Karchmer--Wigderson theorem for De~Morgan formulas, for
protocols specialized to Karchmer--Wigderson games
(Definition~\ref{def:frontiers-kw-protocol} and
Theorem~\ref{thm:frontiers-kw-general}; the same result is
Theorem~\ref{thm:lowerbounds-algebraic-kw}). General two-party communication
complexity is not formalized: protocols for arbitrary functions and relations,
randomized protocols, and lower-bound methods such as rectangle partitions,
fooling sets and discrepancy.

\begin{definition}[Deterministic protocols]
  \label{def:frontiers-protocol}
  A deterministic protocol for $f : X \times Y \to Z$ is a binary tree. Each
  internal node is owned by Alice, who branches on a function of $x$, or by Bob,
  who branches on a function of $y$. Each leaf is labelled by an output. The
  cost is the depth, and $D(f)$ is the minimum cost of a protocol computing $f$.
  Only the protocols of Karchmer--Wigderson games are formalized
  (Definition~\ref{def:frontiers-kw-protocol}); their leaves name a
  coordinate, and they solve a relation rather than compute a function.
\end{definition}

\begin{theorem}[Discrepancy bound for inner product]
  \label{thm:frontiers-ip-discrepancy}
  \uses{def:frontiers-protocol}
  \notready
  For every distribution $\mu$ on $X \times Y$ and every error
  $\varepsilon < 1/2$, public-coin randomized protocols satisfy
  $R_\varepsilon(f) \ge \log_2\bigl((1 - 2\varepsilon)/\mathrm{disc}_\mu(f)\bigr)$.
  The inner product $\mathrm{IP}_n(x,y) = \sum_i x_i y_i \bmod 2$ has uniform
  discrepancy at most $2^{-n/2}$. Hence $R_{1/3}(\mathrm{IP}_n) \ge n/2 - O(1)$.
\end{theorem}
\begin{proof}
  \uses{thm:analysis-parity-orthonormal}
  A $c$-bit protocol partitions $X \times Y$ into at most $2^c$ rectangles.
  Lindsey's lemma bounds each rectangle's bias using orthogonality of the
  Hadamard matrix, that is, of the parities. First step: prove the rectangle
  partition lemma and the fooling-set bound $D(\mathrm{EQ}_n) \ge n$.
\end{proof}

\begin{definition}[Karchmer--Wigderson games and their protocols]
  \label{def:frontiers-kw-protocol}
  \lean{Algebraic.KW.Formula, Algebraic.KW.Formula.eval,
    Algebraic.KW.Formula.depth, Algebraic.KW.Formula.leaves,
    Algebraic.KW.Formula.Computes, Algebraic.KW.Protocol,
    Algebraic.KW.Protocol.run, Algebraic.KW.Protocol.depth,
    Algebraic.KW.Protocol.leaves, Algebraic.KW.Protocol.Solves,
    Algebraic.KW.Protocol.SolvesKW, Algebraic.KW.formulaDepth,
    Algebraic.KW.formulaSize, Algebraic.KW.protocolDepth,
    Algebraic.KW.protocolSize}
  \leanok
  \uses{def:bitstring}
  These definitions come from the imported algebraic-circuits library. A
  De~Morgan formula on $n$ inputs is a binary tree whose leaves are literals
  $x_i$ or $\lnot x_i$ or Boolean constants and whose internal nodes are
  binary AND or OR gates. A protocol on $n$ inputs is a binary tree whose
  leaves each name a coordinate $i < n$ and whose internal nodes are each
  owned by Alice or by Bob and branch by an arbitrary Boolean function of the
  owner's whole input in $\{0,1\}^n$; on inputs $x$ for Alice and $y$ for Bob
  it outputs the coordinate at the leaf it reaches. For both, the depth is the
  length of a longest root-to-leaf path (a leaf has depth $0$) and the size is
  the number of leaves, constant leaves included. A protocol solves the game on
  a rectangle $A \times B$ if its output $i$ satisfies $x_i \neq y_i$ for all
  $x \in A$ and $y \in B$; it solves the Karchmer--Wigderson game of
  $f : \{0,1\}^n \to \{0,1\}$ if it solves the game on
  $f^{-1}(1) \times f^{-1}(0)$. The minimum formula depth and size of $f$ (over
  formulas computing $f$ at every input) and the minimum protocol depth and size
  of its game are infima in $\mathbb{N} \cup \{\infty\}$, equal to $\infty$ when
  no formula or protocol exists.
\end{definition}

\begin{theorem}[Karchmer--Wigderson, general form]
  \label{thm:frontiers-kw-general}
  \lean{Algebraic.KW.formulaDepth_eq_protocolDepth,
    Algebraic.KW.formulaSize_eq_protocolSize}
  \leanok
  \uses{def:frontiers-kw-protocol}
  In the model of Definition~\ref{def:frontiers-kw-protocol}, for every
  $n \ge 1$ and every $f : \{0,1\}^n \to \{0,1\}$, constant or not, the minimum
  depth of a De~Morgan formula computing $f$ equals the minimum depth of a
  protocol for the Karchmer--Wigderson game of $f$, in which Alice holds
  $x \in f^{-1}(1)$, Bob holds $y \in f^{-1}(0)$, and they must output some $i$
  with $x_i \neq y_i$. Likewise the minimum number of leaves of a formula
  computing $f$ equals the minimum number of leaves of a protocol for the game.
  The hypothesis $n \ge 1$ cannot be dropped: for $n = 0$ no protocol exists,
  since a leaf must name a coordinate, while a constant formula has depth $0$.
  This is the result of Theorem~\ref{thm:lowerbounds-algebraic-kw}.
\end{theorem}
\begin{proof}
  \leanok
  Both translations preserve depth and leaf count exactly. From a formula: an
  OR gate becomes an Alice node where she says whether her input satisfies the
  left disjunct, and play continues in a disjunct her input satisfies; an AND
  gate becomes a Bob node where he says whether his input satisfies the left
  conjunct, and play continues in a conjunct his input violates. A literal leaf
  names its coordinate, and a constant leaf, reachable only on an empty
  rectangle, names a default coordinate, which exists because $n \ge 1$. From a
  protocol: track the rectangle of inputs consistent with the transcript. Alice
  nodes become OR gates, Bob nodes become AND gates, and a leaf naming $i$
  becomes whichever of $x_i$ and $\lnot x_i$ is $1$ on Alice's side and $0$ on
  Bob's side, or a constant when one side is empty.
\end{proof}

\section{Algebraic complexity}

The imported algebraic-circuits library already has arithmetic circuits: CSLib
straight-line circuits over its arithmetic signature
(\texttt{Algebraic.Arithmetic.signature}), with binary addition, binary
multiplication and constant gates, evaluable in any carrier with addition and
multiplication, in particular as formal polynomials in Mathlib's
\texttt{MvPolynomial}. They are used for lower bounds on computing a single
polynomial, such as Theorem~\ref{thm:lowerbounds-algebraic-hessian}. Polynomial
families, degree bounds, projections and the classes below are not formalized.

\begin{definition}[$\cc{VP}$ and $\cc{VNP}$]
  \label{def:frontiers-vp-vnp}
  Fix a field $\mathbb{F}$. An arithmetic circuit is a DAG with $+$ and $\times$
  gates over variables and field constants. $\cc{VP}$ is the class of polynomial
  families $(f_n)$ with $\mathrm{poly}(n)$ variables and degree, computed by
  $\mathrm{poly}(n)$-size arithmetic circuits. $\cc{VNP}$ is the class of families
  $f_n(x) = \sum_{e \in \{0,1\}^{m(n)}} g_n(x, e)$ with $(g_n) \in \cc{VP}$ and
  $m$ polynomial.
\end{definition}

\begin{theorem}[Valiant: completeness of the permanent]
  \label{thm:frontiers-permanent}
  \uses{def:frontiers-vp-vnp}
  \notready
  Over any field of characteristic other than $2$, the permanent family
  $(\mathrm{perm}_n)$ is $\cc{VNP}$-complete under polynomial-size projections.
\end{theorem}
\begin{proof}
  Show the permanent is in $\cc{VNP}$ via Ryser's formula. For hardness, convert
  formulas to weighted graphs whose cycle covers compute the polynomial. First
  step: on top of the imported arithmetic circuits and their
  \texttt{MvPolynomial} semantics, define polynomial families, their size and
  degree bounds, and $p$-projections. The Boolean counterpart,
  $\#\cc{P}$-completeness of the $0/1$ permanent, would connect to the existing
  class $\#\cc{P}$ (Definition~\ref{def:sharpP}).
\end{proof}

\section{Fine-grained and parameterized complexity}

\begin{theorem}[SETH implies quadratic hardness of Orthogonal Vectors]
  \label{thm:frontiers-seth-ov}
  \uses{def:dtime}
  \notready
  Suppose that for every $\delta > 0$ there is $k$ such that $k$-SAT on $n$
  variables has no algorithm running in time $2^{(1-\delta)n}\,\mathrm{poly}(n)$
  (SETH). Then for every $\delta > 0$ there is $c$ such that Orthogonal Vectors on
  $N$ vectors of dimension $c \log N$ has no algorithm running in time
  $N^{2-\delta}$.
\end{theorem}
\begin{proof}
  Split the variables into two halves and enumerate partial assignments. Map each
  to a vector over the clauses, after sparsification. First step: running times
  are measured in a parameter (the number of variables $n$) rather than input
  length, so fine-grained reductions need explicit time bounds in named
  parameters, beyond the existing $\cc{DTIME}$.
\end{proof}

\begin{definition}[$\cc{FPT}$ and parameterized reductions]
  \label{def:frontiers-fpt}
  \uses{def:language, def:dtime}
  A parameterized problem is a language $L$ with a polynomial-time parameter
  $\kappa$. It is in $\cc{FPT}$ if it is decidable in time
  $g(\kappa(x)) \cdot |x|^{O(1)}$ for some computable $g$. An fpt-reduction
  preserves membership, runs in fpt time, and bounds the target parameter by a
  function of the source parameter. $\cc{W[1]}$ is the closure of $k$-Clique
  under fpt-reductions. First targets: Vertex Cover is in $\cc{FPT}$ by a
  $2^k$-branching algorithm, and $k$-Clique is $\cc{W[1]}$-complete by
  definition, with $k$-Independent Set equivalent to it.
\end{definition}

\section{Quantum computation}

\begin{theorem}[Adleman--DeMarrais--Huang: $\cc{BQP} \subseteq \cc{PP}$]
  \label{thm:frontiers-bqp-pp}
  \uses{def:PP}
  \notready
  Define $\cc{BQP}$ by polynomial-time uniform families of quantum circuits over
  the gate set $\{\mathrm{Hadamard}, \mathrm{Toffoli}\}$, with error $1/3$. Then
  $\cc{BQP} \subseteq \cc{PP}$.
\end{theorem}
\begin{proof}
  With Hadamard and Toffoli gates all amplitudes are real: each amplitude is a
  signed count of computation paths divided by $2^{h/2}$, where $h$ is the
  number of Hadamard gates. Squaring, the acceptance probability is
  $(\#\text{positive} - \#\text{negative})/2^h$ over pairs of polynomially
  checkable paths that end in the same accepting basis state, which a $\cc{PP}$
  machine can compare with a threshold. First step: define the path-sum
  semantics of such circuits and prove it agrees with the matrix semantics on
  small examples.
\end{proof}
